\chapter{Continuum Limits and Scaling Windows}
\label{ch:constructive-continuum-limits-scaling-windows}

\ConventionsEuclidean

A lattice or cutoff theory becomes a continuum QFT only after its correlation
functions have limiting distributions with the desired locality, covariance,
positivity, and clustering properties.  The continuum limit is therefore a
statement about a family of regulated theories and a choice of coordinates in
the space of local observables.

\section{Scaling Trajectories}

Let \(a>0\) be a lattice spacing and \(\Lambda_a=a\mathbb Z^D\) the lattice.
For each \(a\), let \(S_a[\phi]\) be the regulated action and
\(\langle\cdots\rangle_a\) the corresponding finite-volume expectation,
with the infinite-volume limit taken before the continuum limit when it
exists.  A scaling trajectory is a map
\[
  a\longmapsto c(a)
\]
from lattice spacing to the couplings and counterterms appearing in
\(S_a\).  The word counterterm here means a local term in the regulated
action chosen as part of the path integral definition.

For a lattice observable \(O^I_a(n)\), choose a renormalization matrix
\(Z^I{}_J(a)\) and define the smeared continuum-coordinate observable
\[
  O^I_a(f)
  =
  a^D\sum_{n\in\mathbb Z^D}
  f(an)\,Z^I{}_J(a)\,O^J_a(n),
  \qquad f\in C_c^\infty(\mathbb R^D).
\]
A continuum limit of the \(k\)-point functions is the distributional limit
\[
  \lim_{a\to0}
  \langle O^{I_1}_a(f_1)\cdots O^{I_k}_a(f_k)\rangle_a
  =
  W^{I_1\ldots I_k}(f_1,\ldots,f_k).
\]

\section{OS Reconstruction Target}

The limiting Euclidean distributions define a Wightman QFT only if they
satisfy the Osterwalder--Schrader hypotheses with the growth assumptions
needed for reconstruction.  Reflection positivity is inherited from the
lattice when the reflection-positive structure is uniform in the limit.
Euclidean invariance is obtained only after proving that lattice rotation
breaking terms vanish in the scaling trajectory.

\begin{definition}[Constructive continuum limit]
\label{def:constructive-continuum-limit}
A constructive continuum limit of a regulated family is a collection of
limiting Schwinger distributions satisfying symmetry, Euclidean covariance,
reflection positivity, permutation symmetry, clustering, and the required
regularity and growth bounds, together with an OS reconstruction of the
Hilbert space, vacuum, fields, and local operator domains.
\end{definition}

\section{Scaling Window and Correlation Length}

Let \(\xi(a)\) be a correlation length measured in lattice units.  A
continuum scaling window requires
\[
  \xi(a)\to\infty,\qquad a\,\xi(a)\to \ell_{\rm phys}
\]
with \(\ell_{\rm phys}\) finite or infinite according to whether the target
continuum theory is massive or scale invariant.  Local lattice observables
probe continuum physics at separations \(r\) satisfying
\[
  a\ll r\ll a\,\xi(a)
\]
in a massive scaling limit.  The lower inequality suppresses lattice
artifacts; the upper inequality keeps the observable inside the universal
scaling regime before exponential decay dominates.

\section{Renormalization Coordinates}

Let \(\mathcal B\) be a Banach space of local lattice interactions near a
candidate fixed point \(S_*\), and let \(\mathcal R_b\) be a block-spin map
with scale factor \(b>1\).  Linearization gives
\[
  D\mathcal R_b|_{S_*}\,v_i=b^{y_i}v_i .
\]
Coordinates with \(y_i>0\) are relevant, \(y_i=0\) marginal, and \(y_i<0\)
irrelevant in this specified norm and blocking scheme.  A statement that a
scaling trajectory reaches \(S_*\) is a theorem only after the nonlinear map
\(\mathcal R_b\), its domain, and estimates controlling irrelevant directions
are supplied.

\section{Status of Scalar Examples}

Superrenormalizable scalar theories such as \(P(\phi)_2\) and
\(\phi^4_3\) have constructive continuum realizations with Wightman or OS
reconstruction under their stated hypotheses.  This is different from proving
that the three-dimensional Ising critical point is a conformal field theory
with the full expected operator spectrum and OPE.  In four dimensions, the
existence of a nonperturbative ultraviolet-complete QFT agreeing with
\(\phi^4_4\) to all orders in perturbation theory is not known; perturbative
renormalizability and nonperturbative continuum existence are distinct
claims.

\begin{openproblem}[Nonperturbative fixed points from Wilsonian maps]
\label{op:nonperturbative-wilsonian-fixed-points}
Develop a Wilsonian Banach-space framework for physically relevant
four-dimensional and three-dimensional QFTs in which fixed points,
renormalized operator coordinates, contact terms, and OS reconstruction can
all be controlled in the same theorem.  The framework should make precise
when a physicists' fixed-point argument is a proof and when it is a
well-supported construction principle.
\end{openproblem}
